\subsection{Scaling in Theory} \label{sub:ScalingInTheory}
This section conducts a theoretical analysis of how to arbitrarily scale \textit{DUORAM} to fulfil the objective of this research.
It is based on reading and writing protocol from \ref{sub:AdvancedReadingAndWriting} taken from Vadapalli et al. \cite{vadapalli2023duoram}.

\subsubsection*{General Assumptions} \label{subsub:ScalingInTheoryGeneralAssumptions}
This paragraph highlights general thoughts and assumptions on arbitrarily scaling duoram entries.
The affected values are independent of the phase and operation executed, thus placed separately in this subsection.

The theoretical protocol is not bound to computational limits.
Therefore, the database $D$ is unlimited in the number of entries it can hold, solely the entry size is fixed.
This means, scaling the database equals scaling the entry size of every entry $D[i]$.

As both $D$ and its blinded equivalent $\widetilde{D}$ with its respective additive shares $\widetilde{D}_0,\widetilde{D}_1$ are of the same size, the blinded shares will automatically adjust their size once $D$ is scaled.
This will also affect its respective shares $D_0,D_1$ which are the same size as $D$.
Since $\widetilde{D}_b\leftarrow D_b+\xi_b\text{ for } b\in\{0,1\}$ the blinding factors $\xi_0,\xi_1$ will grow in their entry size too.
Note that all instances of $D,D_0,D_1,\xi_0,\xi_1,\widetilde{D}_0,\widetilde{D}_1$ are of equal size.
They therefore grow at the same rate when the size of $D$ is adjusted keeping the same size.

These general assumptions are taken as granted in this theoretic part.

\subsubsection*{Scaling during Preprocessing} \label{subsub:ScalingDuringPreprocessing}
As preprocessing for reading and writing is nearly similar in their goals and processes, both preprocessing phases are covered in this subsection.
As the main goal of preprocessing is to generate RDPF-Triples for each peer, this process is affected the most during scaling and will affect its derived values.

\paragraph{Scaling DPF Generation}\label{par:ScalingInTheoryDPFGeneration}
The theoretic DPF generation process follows the introduced construction of Boyle et al. \cite{boyle2016function} from \ref{sub:BGIDPFConstruction}.
It takes an $n$-bit seed as input and recursively uses a length-doubling PRG to generate two binary trees.
With correction words, flag bits, and tree invariants both trees are adjusted to represent a DPF.

Vadapalli et al. \cite{vadapalli2023duoram} extend this concept by further using the results of this DPF construction.
Each peer is assigned three DPF shares in form of RDPF-Triples.
These are then evaluated leaving each peer with three vector pairs $(v_b^k, t_b^k)$ with the known properties.

To reconstruct the effects of adjusting the database entry-size the author chose to traverse the algorithm of the BGI-construction (\ref{sub:BGIDPFConstruction}) in reverse.
Recall the properties of $v_b^k$ representing $v^k = e_{i^*}\cdot M$ when combined.
This means $v^k$ is a vector holding $M$ at the target index $i^*$ and zero everywhere else.
When additively shared among both peers, it is divided into two vectors $v_0^k,v_1^k$ holding additive shares of $M$ at the target index and random values everywhere else.
These random values will add to $0$ once combined.
As each entry has the same size and $v_b^k[i^*]$ has size $|M|$ each entry of $v_b^k$ has the same size as the message.
As $D[i]$ is adjusted in its size, $M$ will adjust in the same manner allowing the client to write messages of the new size.
This means each entry $v_b^k[i]$, having the same size as $M$, will also adjust its size once $D[i]$ has changed.

Recall $v_b^k[i]$ being the leaf label at index $i$ from $P_b$'s $k$'th DPF-tree.
As $v_b^k$ adjusts to the size of $M$, each leaf label will have the size of $M$.
Since the BGI-construction uses a length-doubling PRG, the seed of the tree will have size $|M|$ too.
A single execution of the length-doubling PRG therefore generates two children of size $|M|$ each.
This means a single length-doubling PRG execution on a seed of size $|M|$ results in a string of size $2|M|$

As mentioned previously, $F^k_b$ is the negative sum of the according leaf labels $-\sum^n_i{v_b^k[i]}$.
Thus, $F$ will also adjust once $D[i]$ has adjusted its size.

$t_b^k$ represents the flag bits of each leaf in $P_b$'s $k$'th DPF-tree.
As each leaf holds a single flag bit $t_b^k$'s size is solely dependent from the amount of leafs in the DPF.
This means the size of $t_b^k$ will not be affected by scaling $D[i]$.

To conclude, adjusting the entry size of $D[i]$ results in $v_b^k,F^k_b,\text{ and }M_b$ adjusting their size in the same manner.
Changing the size of these values will result in size-changes while generating DPFs and in the preprocessing phase in general.

\subsubsection*{Online-Phase Scaling} \label{subsub:ScalingDuringOnlinePhase}
The online-phase does not yield any new values but the indices $ri$ and $i$.
They indicate which element of the database to retrieve, thus affecting the offset $S$ during cyclic shifting.
This, however, is not dependent from the size of entry $D[b]$ but the amount of entries.
If the datatype for the indices $i$ and $ri$ is large enough to uniquely identify each leaf, no changes are necessary. 